\documentclass[multi=tikzpicture,border=7pt]{standalone}
\usepackage{ctex,amsmath,bm}
\usepackage{tikz}
\usetikzlibrary{arrows.meta,calc,angles,quotes,decorations.pathmorphing,patterns,positioning,intersections,3d}
\tikzset{
  axis/.style={->,>=Stealth,thick},
  force/.style={->,>=Stealth,very thick,red!72!black},
  flow/.style={->,>=Stealth,thick,blue!70!black},
  dim/.style={<->,>=Stealth,thin},
  body/.style={draw=black,fill=blue!7,thick},
  guide/.style={text=black,dashed,gray!70,text=black},
  hidden/.style={dashed,gray!62,thick},
  point/.style={circle,fill=black,inner sep=1.5pt},
  field/.style={->,>=Stealth,thick,purple!75!black},
  panel/.style={font=\bfseries,anchor=west},
  callout/.style={draw=gray!65,rounded corners,fill=white,align=center,inner sep=3pt}
}
\begin{document}

% 01 question-01
\begin{tikzpicture}[scale=2]
 \coordinate (C) at (0,0);\coordinate (O) at (.5,.866025);
 \coordinate (A) at (.697465,.323494);\coordinate (B) at (1.05522,-.60923);
 \draw[very thick] (-1.4,.80829)--(2.2,-1.27017);
 \draw[body] (O) circle(1);
 \draw[brown!70!black,very thick] (A)--(B);
 \draw[blue!70!black,thick] (O)--(A);
 \foreach \p in {O,A,B,C}{\fill (\p) circle(.8pt);}
 \node[above left] at (O){$O$};\node[left] at (A){$A$};\node[right] at (B){$B$};\node[below left] at (C){$C$};
 \draw[guide] (O)--++(0,-1.5);
 \draw[dim] (O)--(C) node[midway,left]{$r$};
 \node[right] at (.75,.67){$OA=r/\sqrt3$};\node[right] at (1.05,-.12){$AB=r$};
 \draw[->] (.5,.446025) arc(-90:-70:.42) node[below right]{$\varphi$};
 \draw[thick] (1.3,-1.27017)--(2.2,-1.27017);
 \draw (1.82,-1.27017) arc(180:150:.38) node[left]{$30^\circ$};
 \node[align=center] at (.4,2.3){圆筒与斜面相切，杆端 $B$ 位于斜面\\$\mu=1/\sqrt3$，橡胶垫摩擦因数为 $\mu'$};
\end{tikzpicture}

% 02 solution-01
\begin{tikzpicture}[scale=1.6]
 \begin{scope}[shift={(-2.2,0)}]
  \coordinate (C) at (0,0);\coordinate (O) at (.5,.866025);
  \draw[very thick] (-1.3,.750555)--(1.8,-1.03923);
  \draw[body] (O) circle(1);
  \draw[guide] (C)--(60:2.1) node[above]{$N$ 方向};
  \draw[guide] (C)--(30:1.8) (C)--(90:1.8);
  \draw[force] (C)--(75:1.4) node[left]{$\bm R_C$};
  \draw (60:.65) arc(60:90:.65) node[above]{$\alpha=30^\circ$};
  \fill (O) circle(1pt) node[right]{$O$};\node[below left] at (C){$C$};
  \node[align=center] at (.3,2.5){(a) 摩擦锥绕斜面法线展开\\边界方向为 $30^\circ$ 与 $90^\circ$};
 \end{scope}
 \begin{scope}[shift={(2.3,0)}]
  \coordinate (O) at (.5,.866025);\coordinate (A) at (.919259,.469286);\coordinate (B) at (.919259,-.530714);
  \draw[very thick] (-.9,.519615)--(1.8,-1.03923);
  \draw[blue!70!black,very thick] (O)--(A);\draw[brown!70!black,very thick] (A)--(B);
  \draw[guide] (O)--++(0,-1.6);\draw[guide] (O)--++(-.7,-1.212436);
  \foreach \p in {O,A,B}{\fill (\p) circle(1pt) node[right]{$\p$};}
  \draw[dim] ($(A)+(.3,0)$)--($(B)+(.3,0)$) node[midway,right]{$r$};
  \node[left] at (.6,.68){$r/\sqrt3$};
  \draw (.5,.45) arc(-90:-43.41784:.416025) node[right]{$\gamma$};
  \node[align=center] at (.4,2.5){(b) 极限构型：杆轴竖直\\投影沿斜面法线取值};
  \node[align=center] at (.4,-1.65){$\cos(\gamma+30^\circ)=\sqrt3-3/2$\\$\gamma\simeq46.58^\circ$};
 \end{scope}
\end{tikzpicture}

% 03 question-02
\begin{tikzpicture}[x={(1cm,0cm)},y={(.56cm,.34cm)},z={(0cm,1cm)}]
  \node[panel] at (-5.1,0,4.8){空间铰接杆系（透视图）};
  \coordinate (O) at (0,0,4.0);
  \coordinate (A) at (0,0,3);
  \coordinate (B) at (.5,0,2.133975);
  \coordinate (D) at (.643923,-1.29212,1.59688);
  \coordinate (E) at (.643923,1.29212,1.59688);
  \coordinate (C) at (.787846,0,1.059785);
  \draw[axis] (-3.4,0,0)--(3.5,0,0) node[right]{$x$};
  \draw[axis] (0,-2.4,0)--(0,2.8,0) node[above right]{$y$};
  \draw[axis] (0,0,-1)--(0,0,4.6) node[above]{$z$};
  \fill[blue!5,opacity=.5] (D)--(B)--(E)--(C)--cycle;
  \draw[hidden] (B)--(E) (E)--(O);
  \draw[guide] (O)--(A);\draw[black,very thick] (A)--(B)--(D)--(O);
  \draw[black,very thick] (D)--(C)--(E);
  \draw[guide] (B)--(C);
  \foreach \P in {O,A,B,C,D,E}{\shade[ball color=orange!48] (\P) circle(.13);}
  \node[above] at (O){$O$}; \node[left] at (A){$A$}; \node[above right] at (B){$B$};
  \node[below right,xshift=2pt] at (C){$C$}; \node[left] at (D){$D$}; \node[right] at (E){$E$};
  \draw[force] (C)--++(0,0,-1.15) node[midway,right]{$m\bm g$};
  \node[callout] at (4.2,-.8,3.2){$OA=AB=L$\\$BD=BE=DC=EC=a$\\$OD=OE=b$};
  \node[font=\small,anchor=west] at (-4.2,2.7,-1.25){粗虚线：后侧杆；细虚线：辅助线};
  \node[font=\small,anchor=west] at (-4.2,2.7,-1.70){$D,E$：空间横向节点};
  \node[font=\small,anchor=west] at (-4.2,2.7,-2.15){$C$：前层节点};
\end{tikzpicture}

% 04 solution-02
\begin{tikzpicture}[x={(1cm,0cm)},y={(.58cm,.34cm)},z={(0cm,1cm)}]
  \begin{scope}[shift={(-3.8,0,0)}]
    \node[panel] at (-2.8,0,4.25){(a) 对称坐标 $DE=2x,\ BC=2y$};
    \coordinate (M) at (0,0,1.5);
    \coordinate (D) at (-2.15,-1.05,1.45);
    \coordinate (E) at (2.15,-1.05,1.45);
    \coordinate (B) at (0,0,2.85);
    \coordinate (C) at (0,2.75,-.55);
    \fill[blue!7,opacity=.65] (D)--(B)--(E)--(C)--cycle;
    \draw[hidden] (B)--(E);
    \draw[very thick] (D)--(B)--(E)--(C)--(D) (B)--(C);
    \draw[dim] (D)--(E) node[midway,below]{$2x$};
    \draw[dim] (B)--(C) node[midway,right]{$2y$};
    \draw[guide] (B)--(M)--(C);
    \fill (B) circle(2pt) node[above]{$B$};
    \fill (C) circle(2pt) node[below right]{$C$};
    \fill (D) circle(2pt) node[left]{$D$};
    \fill (E) circle(2pt) node[right]{$E$};
    \node[font=\small,align=center] at (0,2.8,-1.35){$D,E$ 沿空间横向分开，$B,C$ 位于中垂面};
    \node[callout] at (0,-1.2,-1.0){$x^2+y^2=a^2$};
  \end{scope}
  \begin{scope}[shift={(4.1,0,0)}]
    \node[panel] at (-2.7,0,4.25){(b) 小振动广义坐标 $\theta$};
    \draw[guide] (0,0,-.7)--(0,0,4.0);
    \coordinate (O) at (0,0,3.55);
    \coordinate (A) at (.55,0,2.72);
    \coordinate (B) at (1.10,0,1.89);
    \draw[very thick] (O)--(A)--(B);
    \fill (O) circle(2pt) node[above]{$O$};
    \fill (A) circle(2pt) node[left]{$A$};
    \fill (B) circle(2pt) node[right]{$B,m$};
    \coordinate (AxisBelow) at (0,0,1.4);
    \draw pic[draw,"$\theta$",angle radius=.6cm,angle eccentricity=1.45]{angle=AxisBelow--O--A};
    \draw[dim] (O)--(A) node[midway,left]{$L$};
    \draw[dim] (A)--(B) node[midway,right]{$L$};
    \node[callout] at (0,-.8,.1){$\displaystyle I_{\rm eff}=m\!\left[L^2+\left(\frac{b^2-a^2}{4L}\right)^2\right]$};
  \end{scope}
\end{tikzpicture}

% 05 question-03
\begin{tikzpicture}[x=.92cm,y=.92cm]
  \node[panel] at (-5.0,4.1){同轴圆柱壳与固定长螺线管（横截面）};
  \fill[blue!5] (0,0) circle(3.25);
  \draw[very thick] (0,0) circle(3.25);
  \draw[very thick] (0,0) circle(2.35);
  \draw[very thick] (0,0) circle(1.35);
  \draw[purple!72!black,double,very thick] (0,0) circle(2.35);
  \draw[dim] (0,0)--(1.35,0) node[midway,below]{$a$};
  \draw[dim] (0,0)--({2.35*cos(135)},{2.35*sin(135)}) node[midway,above left]{$b$};
  \draw[dim] (0,0)--({3.25*cos(65)},{3.25*sin(65)}) node[midway,left]{$c$};
  \node at (0,.55){内壳 $+Q,m$};
  \node at (0,1.82){固定螺线管};
  \node at (0,2.82){外壳 $-Q,M$};
  \foreach \x/\y in {-1.6/0,1.6/0,0/-1.6,0/1.6}{\node[purple!75!black] at (\x,\y){$\odot$};}
  \draw[red!72!black,very thick,->] (-1.7,1.85) arc[start angle=132,end angle=205,radius=2.5];
  \node[red!72!black] at (-3.4,2.5){螺线管电流 $I\downarrow0$};
  \node[callout] at (7,.3){螺线管单位长度匝数 $n$\\$+z$ 入纸，初始 $B$ 出纸\\磁场仅在 $r<b$ 内近似均匀\\内、外壳可绕轴转动；中间为固定螺线管};
\end{tikzpicture}

% 06 solution-03
\begin{tikzpicture}[x=.9cm,y=.9cm]
  \begin{scope}[shift={(-4.15,0)}]
    \node[panel] at (-3.1,3.9){(a) 感生电场与机械角动量};
    \draw[very thick] (0,0) circle(1.2) (0,0) circle(2.2) (0,0) circle(3.1);
    \draw[purple!70!black,dashed,thick] (0,0) circle(2.2);
    \draw[field] (1.2,0) arc[start angle=0,end angle=55,radius=1.2];
    \draw[field] (3.1,0) arc[start angle=0,end angle=55,radius=3.1];
    \node at (1.25,.65){$\bm E_\varphi(a)$};
    \node at (3.2,1.7){$\bm E_\varphi(c)$};
    \node[callout] at (0,-3.8){$E_\varphi(a)=\dfrac{\mu_0na}{2}\dot I$\qquad$E_\varphi(c)=\dfrac{\mu_0nb^2}{2c}\dot I$};
  \end{scope}
  \begin{scope}[shift={(4.25,.05)}]
    \node[panel] at (-3.0,3.85){(b) 场角动量所在环域};
    \fill[blue!10] (0,0) circle(2.75);
    \fill[white] (0,0) circle(1.25);
    \draw[very thick] (0,0) circle(1.25) (0,0) circle(2.75);
    \draw[field] (1.7,0)--(2.5,0) node[right]{$\bm E$};
    \node[purple!75!black] at (0,1.9){$\odot\ \bm B$};
    \draw[flow] (0,-1.85) arc[start angle=-90,end angle=-170,radius=1.85];
    \node[blue!70!black] at (2.3,-1.5){$\bm g=\varepsilon_0\bm E\times\bm B$};
    \node[callout] at (0,-3.55){$\displaystyle L_{\rm em}=\int_a^b r\,\varepsilon_0EB\,(2\pi rl\,dr)$};
  \end{scope}
\end{tikzpicture}

% 07 question-04
\begin{tikzpicture}[x={(1cm,0cm)},y={(.58cm,.35cm)},z={(0cm,1cm)}]
  \begin{scope}[shift={(-4.25,0,0)}]
    \node[panel] at (-2.7,0,4.3){(a) 三维立方体电阻网络};
    \coordinate (p000) at (0,0,0); \coordinate (p100) at (2.6,0,0);
    \coordinate (p110) at (2.6,2.6,0); \coordinate (p010) at (0,2.6,0);
    \coordinate (p001) at (0,0,2.6); \coordinate (p101) at (2.6,0,2.6);
    \coordinate (p111) at (2.6,2.6,2.6); \coordinate (p011) at (0,2.6,2.6);
    \draw[hidden] (p000)--(p010) (p010)--(p110) (p010)--(p011);
    \draw[very thick] (p000)--(p100)--(p110) (p100)--(p101) (p000)--(p001)--(p101)--(p111)--(p110) (p001)--(p011)--(p111) (p011)--(p010);
    \foreach \P/\lab/\pos in {p000/$(000)$/below left,p100/$(100)$/below,p110/$(110)$/right,p010/$(010)$/above left,p001/$(001)$/left,p101/$(101)$/below right,p111/$(111)$/above right,p011/$(011)$/above}{\fill (\P) circle(2pt);\node[\pos,font=\small] at (\P){\lab};}
    \draw[red!72!black,very thick] (p000) circle(.16) (p101) circle(.16);
    \node[callout] at (1.3,-1.2,-.8){每条棱电阻均为 $r$};
  \end{scope}
  \begin{scope}[shift={(4.15,0,0)},scale=.82]
    \node[panel] at (-3.3,0,5.15){(b) 四维超立方体的三维投影};
    \coordinate (s) at (1.0,.7,.8);
    \coordinate (a000) at (0,0,0); \coordinate (a100) at (3,0,0);
    \coordinate (a110) at (3,3,0); \coordinate (a010) at (0,3,0);
    \coordinate (a001) at (0,0,3); \coordinate (a101) at (3,0,3);
    \coordinate (a111) at (3,3,3); \coordinate (a011) at (0,3,3);
    \foreach \P in {000,100,110,010,001,101,111,011}{\coordinate (b\P) at ($(a\P)+(s)$);}
    \foreach \prefix/\sty in {a/black,b/blue!68!black}{
      \draw[\sty,thick] (\prefix000)--(\prefix100)--(\prefix110)--(\prefix010)--cycle
        (\prefix001)--(\prefix101)--(\prefix111)--(\prefix011)--cycle
        (\prefix000)--(\prefix001) (\prefix100)--(\prefix101) (\prefix110)--(\prefix111) (\prefix010)--(\prefix011);}
    \foreach \P in {000,100,110,010,001,101,111,011}{\draw[orange!75!black,thick] (a\P)--(b\P);}
    \fill (a000) circle(2.6pt) node[below left]{$(0000)$};
    \fill[blue!70!black] (b111) circle(2.6pt) node[above right]{$(1111)$};
    \node[callout] at (2.0,-1.5,-.8){黑、蓝两立方体表示第四位 $0,1$\\橙色连线为第四维棱};
  \end{scope}
\end{tikzpicture}

% 08 solution-04
\begin{tikzpicture}[x=.9cm,y=.9cm]
  \begin{scope}[shift={(-4.3,0)}]
    \node[panel] at (-3.1,3.5){(a) 立方体面对角端点的等势缩并};
    \coordinate (A0) at (-2.7,.8); \coordinate (B0) at (2.7,.8);
    \coordinate (C0) at (0,2.35); \coordinate (D0) at (-1.75,-1.15);
    \coordinate (E0) at (1.75,-1.15); \coordinate (F0) at (0,-2.25);
    \draw[very thick] (A0)--(C0)--(B0) (A0)--(D0)--(F0)--(E0)--(B0) (C0)--(F0);
    \fill (A0) circle(2.2pt) node[above]{$A$};
    \fill (B0) circle(2.2pt) node[above]{$B$};
    \fill (C0) circle(2.2pt) node[above]{$\{001,100\}$};
    \fill (D0) circle(2.2pt) node[above]{$010$};
    \fill (E0) circle(2.2pt) node[above]{$111$};
    \fill (F0) circle(2.2pt) node[below]{$\{011,110\}$};
    \node[font=\small] at (-1.4,1.72){$r/2$}; \node[font=\small] at (1.4,1.72){$r/2$};
    \node[font=\small] at (-2.45,-.25){$r$}; \node[font=\small] at (2.45,-.25){$r$};
    \node[font=\small] at (0.35,.15){$r/2$};
    \node[callout] at (0,-3.35){由对称支路并联得 $R_{000,101}=3r/4$};
  \end{scope}
  \begin{scope}[shift={(4.15,.1)}]
    \node[panel] at (-3.1,3.4){(b) 四维对角点的汉明层};
    \foreach \x/\lab/\n in {-2.8/0/1,-1.4/1/4,0/2/6,1.4/3/4,2.8/4/1}{
      \fill[blue!65!black] (\x,1.0) circle(.10);
      \node[font=\small] at (\x,1.52){第 $\lab$ 层};
      \node[font=\small] at (\x,.42){$\n$ 个点};}
    \draw[very thick] (-2.8,1.0)--(-1.4,1.0)--(0,1.0)--(1.4,1.0)--(2.8,1.0);
    \node[fill=white,inner sep=1pt] at (-2.1,-.18){$r/4$};
    \node[fill=white,inner sep=1pt] at (-.7,-.18){$r/12$};
    \node[fill=white,inner sep=1pt] at (.7,-.18){$r/12$};
    \node[fill=white,inner sep=1pt] at (2.1,-.18){$r/4$};
    \node[font=\small,text=gray!70!black] at (0,-.72){相邻两层间各边并联};
    \node[callout] at (0,-1.85){$R_{0000,1111}=2r/3$\\$R_{0000,1101}=61r/96$（节点电势法）};
  \end{scope}
\end{tikzpicture}

% 09 question-05
\begin{tikzpicture}[x={(1cm,0cm)},y={(.5cm,.25cm)},z={(0cm,1cm)}]
 \foreach \h in {0,4}{
  \draw[gray,dashed] plot[domain=0:180,samples=70,variable=\t] ({3*cos(\t)},{3*sin(\t)},\h);
  \draw[black,thick] plot[domain=180:360,samples=70,variable=\t] ({3*cos(\t)},{3*sin(\t)},\h);
 }
 \draw[thick] (-3,0,0)--(-3,0,4) (3,0,0)--(3,0,4);
 \draw[axis] (0,0,-.5)--(0,0,4.7) node[above]{$z$};
 \coordinate (P) at (3,0,.7);\coordinate (G) at (2.4,0,1.8);
 \draw[brown!75!black,very thick] (P)--(1.8,0,2.9);
 \fill (P) circle(2pt) node[right]{$P$};\fill (G) circle(2pt) node[left]{$G$};
 \draw[force] (G)--++(0,0,-1) node[left]{$mg$};
 \draw[dim] (0,0,.7)--(P) node[midway,below]{$R$};
 \draw[flow,blue] plot[domain=310:350,samples=30,variable=\t] ({3.15*cos(\t)},{3.15*sin(\t)},2.3) node[right]{$\omega$};
 \node[align=center] at (0,0,5.5){竖直圆筒内壁的水平圆周运动\\圆筒轴与重力方向平行};
 \node[anchor=west] at (3.3,0,2.7){骑手与车架总长 $h$};
\end{tikzpicture}

% 10 solution-05
\begin{tikzpicture}[x=.88cm,y=.88cm]
  \begin{scope}[shift={(-4.15,0)}]

 \node[panel] at (-2.5,3.8){(a) 径向与竖直方向的受力};
 \draw[very thick] (1.5,-2.5)--(1.5,2.8);
 \coordinate (P) at (1.5,0);\filldraw[fill=orange!35] (P) circle(.15) node[right]{$P$};
 \draw[force] (P)--(-.3,0) node[left]{$N$};
 \draw[force] (P)--++(0,1.7) node[left]{$f$};\draw[force] (P)--++(0,-1.7) node[left]{$mg$};
 \node[left] at (1.5,2.5){竖直筒壁};
 \node[callout] at (0,-3.55){$N=m\omega_0^2R,\quad f=mg$};
 \end{scope}
  \begin{scope}[shift={(4.2,.05)}]
    \node[panel] at (-3.1,3.75){(b) 有限长度杆的力矩};
    \draw[very thick] (-2.9,-1.6)--(2.9,-1.6);
    \coordinate (P) at (1.9,-1.6);
    \coordinate (G) at (.55,.55);
    \draw[brown!75!black,very thick] (P)--(-.55,2.25);
    \fill (P) circle(2pt) node[below right]{$P$};
    \fill (G) circle(2pt) node[above right]{$G$};
    \draw[force] (G)--++(0,-1.7) node[below left,xshift=-2pt]{$mg$};
    \draw[force,purple!75!black] (G)--++(1.65,0) node[right]{$m\omega^2\rho_G$};
    \coordinate (WallLeft) at (-1,-1.6); \coordinate (RodTop) at (-.55,2.25);
    \draw pic[draw,"$\theta$",angle radius=.75cm,angle eccentricity=1.35]{angle=RodTop--P--WallLeft};
    \draw[dim] ($(P)+(-.38,-.24)$)--($(-.55,2.25)+(-.38,-.24)$) node[midway,left]{$h$};
    \node[callout] at (0,-3.25){杆元半径 $\rho(x)=R-x\cos\theta$\\力臂为 $x\sin\theta$};
  \end{scope}
\end{tikzpicture}

% 11 question-06
\begin{tikzpicture}[x=.9cm,y=.9cm]
  \node[panel] at (-5.0,4.2){抛物线物点经薄透镜后的像点轨迹};
  \draw[axis] (-5,0)--(5.2,0) node[right]{$x$};
  \draw[axis] (0,-2.8)--(0,4.2) node[above]{$y$};
  \draw[green!60!black,very thick,domain=-2.5:2.5,samples=100] plot (\x,{.38*\x*\x});
  \node[green!60!black,anchor=east] at (-2.7,3.55){物点轨迹 $y=ax^2$};
  \draw[blue!70!black,very thick,domain=-2.5:2.5,samples=100] plot (\x,{3.15-.34*\x*\x});
  \fill[blue!70!black] (1.6,{3.15-.34*1.6*1.6}) circle(2pt);
  \node[blue!70!black,anchor=west] at (2.15,2.75){$A(x_1,y_1)$};
  \draw[blue!55!black,thin] (1.65,2.32)--(2.08,2.65);
  \fill[blue!70!black] (-2.2,{3.15-.34*2.2*2.2}) circle(2pt);
  \node[blue!70!black,anchor=east] at (-2.55,1.02){$B(x_2,y_2)$};
  \draw[blue!55!black,thin] (-2.23,1.48)--(-2.48,1.12);
  \node[blue!70!black,anchor=west] at (2.75,3.45){像点轨迹};
  \draw[purple!72!black,very thick] (-2.6,1.55) .. controls (-1.1,2.15) and (1.1,.95) .. (2.6,1.55);
  \draw[purple!72!black,very thick] (-2.6,1.55) .. controls (-1.1,.95) and (1.1,2.15) .. (2.6,1.55);
  \fill (0,1.55) circle(2pt) node[above right,yshift=2pt]{$P(0,p)$};
  \node[purple!72!black,anchor=west] at (2.8,.78){薄透镜};
  \draw[dim] (-.55,0)--(-.55,1.55) node[midway,left]{$p$};
  \draw[guide] (-4,3.15)--(4,3.15);
  \node[callout] at (0,-2.2){待求：像抛物线顶点纵坐标 $h$、光心纵坐标 $p$、焦距 $f$};
\end{tikzpicture}

% 12 solution-06
\begin{tikzpicture}[x=.88cm,y=.88cm]
  \begin{scope}[shift={(-4.2,0)}]
    \node[panel] at (-3.1,3.7){(a) 单个物点的成像几何};
    \draw[axis] (-3.3,0)--(3.5,0) node[right]{$x$};
    \draw[axis] (0,-2.2)--(0,3.5) node[above]{$y$};
    \draw[purple!72!black,very thick] (-2.5,1.1)--(2.5,1.1);
    \fill (0,1.1) circle(2pt) node[right]{$P$};
    \coordinate (Q) at (1.65,.20);
    \coordinate (Qi) at (-1.2,1.754545);
    \fill (Q) circle(2pt) node[right]{$Q(x,y)$};
    \fill (Qi) circle(2pt) node[left]{$Q'(X,Y)$};
    \draw[blue!70!black,thick] (Q)--(0,1.1)--(Qi);
    \draw[blue!70!black,thick] (Q)--(1.65,1.1)--(Qi);
    \node[callout] at (0,-1.8){$s=p-y,\quad s'=Y-p$\\$1/f=1/s+1/s'$};
  \end{scope}
  \begin{scope}[shift={(4.25,0)}]
    \node[panel] at (-3.0,3.7){(b) $p=f$ 时的轨迹变换};
    \draw[axis] (-3.2,0)--(3.4,0) node[right]{$X$};
    \draw[axis] (0,-2.1)--(0,3.45) node[above]{$Y$};
    \draw[blue!70!black,very thick,domain=-2.55:2.55,samples=100] plot (\x,{2.75-.38*\x*\x});
    \draw[guide] (-3,2.75)--(3,2.75);
    \fill (0,2.75) circle(2pt) node[above right]{$h=2f$};
    \node[callout] at (0,-1.65){$X=f/(ax),\quad Y=2f-aX^2$};
  \end{scope}
\end{tikzpicture}

% 13 question-07
\begin{tikzpicture}[x=.88cm,y=.88cm]
  \begin{scope}[shift={(-4.25,0)}]
    \node[panel] at (-3.2,3.85){(a) $S$ 系：下层介质向右运动};
    \fill[blue!6] (-3.2,0) rectangle (3.2,3.0);
    \fill[green!7] (-3.2,-2.8) rectangle (3.2,0);
    \draw[very thick] (-3.2,0)--(3.2,0) node[right]{$x$};
    \draw[axis] (0,-2.8)--(0,3.25) node[above]{$y$};
    \node at (-2.5,2.5){$n_1$ 静止}; \node at (-2.5,-2.3){$n_2$};
    \draw[flow] (-1.6,-1.9)--(1.0,-1.9) node[midway,above]{$u$};
    \draw[blue!70!black,very thick,->] (-2.2,2.75)--(0,0);
    \draw[guide] (-2.6,1.55)--(.8,-1.25);
    \coordinate (NormalRay) at (-.55,.7); \coordinate (OriginA) at (0,0); \coordinate (IncidentStart) at (-2.2,2.75);
    \draw pic[draw,"$\theta$",angle radius=.65cm,angle eccentricity=1.45]{angle=NormalRay--OriginA--IncidentStart};
    \fill (0,0) circle(2pt) node[above right]{$A$};
  \end{scope}
  \begin{scope}[shift={(4.25,0)}]
    \node[panel] at (-3.2,3.85){(b) $S'$ 系：惠更斯作图};
    \fill[blue!6] (-3.2,0) rectangle (3.2,3.0);
    \fill[green!7] (-3.2,-2.8) rectangle (3.2,0);
    \draw[very thick] (-3.2,0)--(3.2,0);
    \draw[axis] (0,-2.8)--(0,3.25) node[above]{$y'$};
    \fill (0,0) circle(2pt) node[above left]{$A$};
    \fill (2.5,0) circle(2pt) node[above]{$B$};
    \draw[green!60!black,thick] (0,0) circle(1.4);
    \draw[blue!70!black,very thick] (2.5,0)--({1.4*cos(-50)},{1.4*sin(-50)});
    \draw[guide] (0,0)--({1.4*cos(-50)},{1.4*sin(-50)});
    \coordinate (DownNormal) at (0,-1); \coordinate (OriginB) at (0,0); \coordinate (WavePoint) at ({1.4*cos(-50)},{1.4*sin(-50)});
    \draw pic[draw,"$i'$",angle radius=.65cm,angle eccentricity=1.45]{angle=DownNormal--OriginB--WavePoint};
    \draw[flow] (1.2,2.2)--(-1.1,2.2) node[midway,above]{上层相对速度 $-u$};
  \end{scope}
\end{tikzpicture}

% 14 solution-07
\begin{tikzpicture}[x=.88cm,y=.88cm]
  \begin{scope}[shift={(-4.2,0)}]
    \node[panel] at (-3.2,3.65){(a) 事件间隔的洛伦兹变换};
    \draw[axis] (-.4,0)--(5.6,0) node[right]{$x$};
    \draw[axis] (0,-.3)--(0,4.0) node[above]{$ct$};
    \fill (0,0) circle(2pt) node[below left]{$A$};
    \fill (4.3,2.4) circle(2pt) node[above right]{$B$};
    \draw[blue!70!black,very thick] (0,0)--(4.3,2.4);
    \draw[guide] (4.3,0)--(4.3,2.4);
    \draw[dim] (0,-.5)--(4.3,-.5) node[midway,below]{$\Delta$};
    \draw[dim] (4.75,0)--(4.75,2.4) node[midway,right]{$c\tau$};
    \node[callout] at (2.5,-1.65){$\Delta=c\tau/(n_1\sin\theta)$};
  \end{scope}
  \begin{scope}[shift={(4.25,0)}]
    \node[panel] at (-3.1,3.65){(b) 相邻大气层的掠射条件};
    \draw[very thick] (0,0) circle(2.2);
    \draw[green!60!black,very thick] (0,0) circle(2.75);
    \draw[blue!70!black,very thick,->] (0,2.2)--(.72,2.65);
    \draw[guide] (0,0)--(0,2.75) (0,0)--(.72,2.65);
    \draw[dim] (-.45,0)--(-.45,2.2) node[midway,left]{$R$};
    \draw[dim] (-.85,2.2)--(-.85,2.75) node[midway,left]{$dh$};
    \node[callout] at (0,-3.0){$u=\alpha\,dh$\\$\sin\theta=R/(R+dh)$};
  \end{scope}
\end{tikzpicture}

% 15 question-08
\begin{tikzpicture}[x={(1cm,0cm)},y={(.55cm,.30cm)},z={(0cm,1cm)}]
  \node[panel] at (-5.2,0,4.75){结冰湖面的三维剖切示意};
  \fill[blue!4] (-4,0,0)--(4,0,0)--(4,3,0)--(-4,3,0)--cycle;
  \fill[cyan!14] (-4,0,0)--(4,0,0)--(4,0,-3.4)--(-4,0,-3.4)--cycle;
  \fill[cyan!7] (4,0,0)--(4,3,0)--(4,3,-3.4)--(4,0,-3.4)--cycle;
  \fill[white] (-4,0,.0)--(4,0,.0)--(4,3,.0)--(-4,3,.0)--cycle;
  \fill[blue!20] (-4,0,0)--(4,0,0)--(4,3,0)--(-4,3,0)--cycle;
  \draw[very thick] (-4,0,0)--(4,0,0)--(4,3,0)--(-4,3,0)--cycle;
  \draw[very thick] (-4,0,-3.4)--(4,0,-3.4)--(4,3,-3.4);
  \draw[very thick] (-4,0,0)--(-4,0,-3.4) (4,0,0)--(4,0,-3.4) (4,3,0)--(4,3,-3.4);
  \draw[blue!70!black,thick] (-4,0,-.16)--(4,0,-.16)--(4,3,-.16);
  \draw[flow] (-2.6,1.1,-2.8)--(-2.6,1.1,-.55) node[midway,left]{$q_w$};
  \draw[flow] (2.5,1.2,-.05)--(2.5,1.2,1.75) node[midway,right]{$q_i$};
  \node at (-2.5,2.2,.35){冰层 $d=1\,\mathrm{cm}$};
  \node at (0,2.0,-2.2){水深 $H_0=5\,\mathrm m$};
  \node[anchor=east] at (1.65,2.1,1.75){空气 $-10^\circ\mathrm C$};
  \node at (1.7,2.3,-3.15){湖底水温 $4^\circ\mathrm C$};
\end{tikzpicture}

% 16 solution-08
\begin{tikzpicture}[x=.9cm,y=.9cm]
  \begin{scope}[shift={(-4.15,0)}]
    \node[panel] at (-3.1,3.65){(a) 初始热流密度};
    \draw[very thick] (-2.5,2.3) rectangle (2.5,2.8);
    \fill[blue!18] (-2.5,2.3) rectangle (2.5,2.8);
    \fill[cyan!9] (-2.5,-2.2) rectangle (2.5,2.3);
    \draw[very thick] (-2.5,-2.2) rectangle (2.5,2.3);
    \draw[flow] (-1.0,-1.4)--(-1.0,1.5) node[midway,left]{$q_w=\kappa_w(4/H_0)$};
    \draw[flow] (1.0,2.35)--(1.0,3.55) node[midway,right]{$q_i=\kappa_i(10/d)$};
    \node at (0,-1.75){$4^\circ\mathrm C$}; \node at (0,2.55){$0^\circ\mathrm C$};
  \end{scope}
  \begin{scope}[shift={(4.2,0)}]
    \node[panel] at (-3.1,3.65){(b) 最终稳态与质量守恒};
    \draw[very thick] (-2.4,-2.3) rectangle (2.4,2.8);
    \fill[cyan!9] (-2.4,-2.3) rectangle (2.4,-1.45);
    \fill[blue!18] (-2.4,-1.45) rectangle (2.4,2.8);
    \draw[flow] (0,-1.9)--(0,2.2) node[midway,right]{$q$};
    \draw[dim] (-2.8,-2.3)--(-2.8,-1.45) node[midway,left]{$h_w$};
    \draw[dim] (-2.8,-1.45)--(-2.8,2.8) node[midway,left]{$h_i$};
    \node[callout] at (0,-3.15){$\kappa_i\dfrac{10}{h_i}=\kappa_w\dfrac4{h_w}$\\$\rho_i h_i+\rho_wh_w=\rho_wH_0$};
  \end{scope}
\end{tikzpicture}

\end{document}
